\documentclass[12pt,a4paper]{article}
\usepackage{fontspec}
\usepackage{polyglossia}
\setmainlanguage{vietnamese}
\setmainfont{DejaVu Serif}
\setsansfont{DejaVu Sans}
\setmonofont{DejaVu Sans Mono}
\usepackage{geometry}
\geometry{top=2cm, bottom=2cm, left=2.5cm, right=2cm}
\usepackage{longtable}
\usepackage{array}
\usepackage{amssymb}
\usepackage{indentfirst}
\usepackage{booktabs}
\usepackage{xcolor}
\usepackage{colortbl}
\usepackage{enumitem}

\definecolor{headerblue}{RGB}{220,230,242}
\definecolor{notegreen}{RGB}{234,250,234}

\begin{document}

\begin{center}
    \Large\textbf{XÁC ĐỊNH YÊU CẦU HỆ THỐNG} \\
    \vspace{0.2cm}
    \large\textbf{Đề tài: Hệ thống nhận diện bệnh trên cây lúa dựa trên thuật toán học sâu CNN}
\end{center}

% ============================================================
\section*{Phạm vi áp dụng đề tài}
\begin{itemize}
    \item \textbf{Phạm vi nghiệp vụ:} Hệ thống số hóa quy trình chẩn đoán bệnh lúa qua hình ảnh bằng công nghệ nhận diện chuyên sâu. Hệ thống kết hợp phân tích hình ảnh và các thông số môi trường để đưa ra chẩn đoán chính xác. Kết hợp xây dựng kho tri thức về dinh dưỡng thực vật để gợi ý các loại vitamin và vi lượng cần bổ sung dựa trên việc truy xuất dữ liệu thông minh.
    \item \textbf{Phạm vi đối tượng:} Phục vụ Nông dân (nhận diện \& nhận gợi ý dinh dưỡng) và Quản trị viên (quản lý danh mục, cập nhật mô hình, quản lý kho kiến thức hệ thống).
    \item \textbf{Giới hạn đề tài:} Hệ thống tự chủ kho dữ liệu dinh dưỡng nội bộ. Đối với thông số môi trường, hệ thống sử dụng tọa độ vị trí (GPS) để tự động truy xuất dữ liệu thời tiết làm dữ liệu phụ trợ. Hỗ trợ nhận diện nhiều loại bệnh cùng lúc trên cùng một ảnh.
\end{itemize}

% ============================================================
\section*{Các bước xác định yêu cầu}

Quá trình thực hiện xác định yêu cầu gồm 2 bước chính như sau:
\begin{itemize}
    \item \textbf{Bước 1:} Khảo sát hiện trạng, kết quả nhận được là các báo cáo về hiện trạng chẩn đoán bệnh trên cây lúa và nhu cầu bổ sung dinh dưỡng.
    \item \textbf{Bước 2:} Lập danh sách các yêu cầu, kết quả nhận được là danh sách các yêu cầu sẽ được thực hiện trên hệ thống máy tính.
\end{itemize}

\subsection*{1. Đối tượng tham gia xác định yêu cầu}
Gồm 2 nhóm người:
\begin{itemize}
    \item \textbf{Chuyên viên tin học:} Thiết kế kiến trúc nhận diện thông minh, xây dựng kho kiến thức hệ thống và luồng xử lý truy xuất thông tin dinh dưỡng chính xác nhất.
    \item \textbf{Nhà chuyên môn (Kỹ sư nông nghiệp):} Cung cấp dữ liệu ảnh kèm thông số môi trường tương ứng, gán nhãn bệnh và trực tiếp biên soạn/chuẩn hóa bộ dữ liệu văn bản về vitamin/dinh dưỡng.
\end{itemize}
Hai nhóm người này phối hợp thật chặt chẽ để có thể xác định đầy đủ và chính xác các yêu cầu.

\hrulefill

\subsection*{2. Khảo sát hiện trạng}
\subsubsection*{2.1 Các hình thức thực hiện phổ biến}
\begin{itemize}
    \item \textbf{Quan sát:} Theo dõi và ghi hình quá trình cây lúa biểu hiện triệu chứng thiếu hụt dinh dưỡng hoặc mắc nhiều bệnh cùng lúc trên đồng ruộng dưới các điều kiện thời tiết khác nhau.
    \item \textbf{Khảo sát:} Khảo sát các người nông dân về mối liên hệ giữa các loại bệnh lúa, sự thiếu hụt dinh dưỡng và sự tác động của môi trường (nhiệt độ, độ ẩm).
    \item \textbf{Thu thập thông tin, tài liệu:} Thu thập tập dữ liệu ảnh lá lúa (kèm dữ liệu đặc tả về thời tiết/độ ẩm) và chuẩn bị bộ dữ liệu văn bản về dinh dưỡng thực vật.
\end{itemize}

\subsubsection*{2.2 Quy trình thực hiện}
\textbf{A. Tìm hiểu tổng quan về thế giới thực:}
\begin{itemize}
    \item Quy mô hoạt động nông nghiệp, các bệnh dịch thường gặp.
    \item Xu hướng chuyển dịch sang bổ sung vitamin tăng sức đề kháng thay vì lạm dụng thuốc hóa học.
\end{itemize}

\textbf{B. Tìm hiểu hiện trạng tổ chức (cơ cấu tổ chức):}
\begin{itemize}
    \item Cơ cấu tổ chức bao gồm các chức danh: Quản trị viên (Quản lý dữ liệu bệnh, mô hình, tập dữ liệu dinh dưỡng), Nông dân (tải ảnh và mô tả tình trạng).
\end{itemize}

\textbf{C. Tìm hiểu hiện trạng nghiệp vụ:}
\begin{itemize}
    \item Việc tìm hiểu dựa trên: Thông tin đầu vào, Quá trình xử lý, Thông tin kết xuất.
    \item Xếp loại các nghiệp vụ vào 4 loại: Lưu trữ, Tính toán, Tra cứu, Kết xuất.
\end{itemize}

\hrulefill

\subsection*{3. Lập danh sách các yêu cầu}
\subsubsection*{3.1 Xác định yêu cầu chức năng nghiệp vụ}

% ============================================================
% BỘ PHẬN: NÔNG DÂN
% ============================================================
\section*{A. BỘ PHẬN: NÔNG DÂN (Mã số: ND)}

\textbf{Mục đích:} Sử dụng hệ thống để chẩn đoán bệnh lúa qua hình ảnh, xem gợi ý dinh dưỡng và theo dõi lịch sử chẩn đoán cá nhân.

\vspace{0.3cm}
\textbf{1. Bảng yêu cầu chức năng nghiệp vụ}

\begin{longtable}{|c|p{3cm}|p{2cm}|p{2.5cm}|p{2cm}|p{2.5cm}|}
\hline
\rowcolor{headerblue}
\textbf{STT} & \textbf{Công việc} & \textbf{Loại công việc} & \textbf{Quy định/CT liên quan} & \textbf{Biểu mẫu liên quan} & \textbf{Ghi chú} \\
\hline
\endhead
\multicolumn{6}{|l|}{\cellcolor{notegreen}\textbf{I. Quản lý tài khoản}} \\
\hline
1 & Đăng ký tài khoản & Lưu trữ & ND\_QĐ 1 & ND\_BM 1 & \\
\hline
2 & Đăng nhập hệ thống & Tra cứu & ND\_QĐ 2 & & \\
\hline
3 & Khôi phục mật khẩu & Lưu trữ & ND\_QĐ 3 & & \\
\hline
\multicolumn{6}{|l|}{\cellcolor{notegreen}\textbf{II. Hồ sơ cá nhân}} \\
\hline
4 & Xem thông tin cá nhân & Tra cứu & ND\_QĐ 4 & & \\
\hline
5 & Chỉnh sửa thông tin cá nhân & Lưu trữ & ND\_QĐ 5 & ND\_BM 2 & \\
\hline
\multicolumn{6}{|l|}{\cellcolor{notegreen}\textbf{III. Chẩn đoán bệnh lúa}} \\
\hline
6 & Gửi ảnh lúa để chẩn đoán & Lưu trữ & ND\_QĐ 6 & ND\_BM 3 & \\
\hline
7 & Nhập thông số môi trường bổ sung & Lưu trữ & ND\_QĐ 7 & & Tùy chọn, không bắt buộc \\
\hline
8 & Xem kết quả chẩn đoán & Tra cứu, Tính toán & ND\_QĐ 8 & & \\
\hline
\multicolumn{6}{|l|}{\cellcolor{notegreen}\textbf{IV. Tra cứu \& Phản hồi}} \\
\hline
9 & Tra cứu lịch sử chẩn đoán & Tra cứu & ND\_QĐ 9 & & \\
\hline
10 & Xem chi tiết một lần chẩn đoán & Tra cứu & ND\_QĐ 10 & & \\
\hline
11 & Gửi phản hồi kết quả & Lưu trữ & ND\_QĐ 11 & ND\_BM 4 & \\
\hline
12 & Xem kết quả xử lý phản hồi & Tra cứu & ND\_QĐ 12 & & \\
\hline
\end{longtable}

% ============================================================
\vspace{0.3cm}
\textbf{2. Bảng Quy định / Công thức liên quan --- Nông dân}

\begin{longtable}{|c|p{1.8cm}|p{3.5cm}|p{6.5cm}|p{1cm}|}
\hline
\rowcolor{headerblue}
\textbf{STT} & \textbf{Mã số} & \textbf{Tên Quy định} & \textbf{Mô tả chi tiết} & \textbf{Ghi chú} \\
\hline
\endhead
1 & ND\_QĐ 1 & Quy định đăng ký tài khoản & Chỉ tạo được tài khoản khi: tên đăng nhập chưa tồn tại, email chưa được đăng ký, mật khẩu đủ độ mạnh. & \\
\hline
2 & ND\_QĐ 2 & Quy định đăng nhập & Kiểm tra tên đăng nhập hoặc email có tồn tại không, sau đó kiểm tra mật khẩu có khớp không. Nhập sai 5 lần thì khóa 15 phút. & \\
\hline
3 & ND\_QĐ 3 & Quy định khôi phục mật khẩu & Thực hiện theo 4 bước: nhập email → nhận OTP → nhập mật khẩu mới → xác nhận mật khẩu mới. & \\
\hline
4 & ND\_QĐ 4 & Quy định xem thông tin cá nhân & Chỉ xem được thông tin của chính mình sau khi đã đăng nhập. & \\
\hline
5 & ND\_QĐ 5 & Quy định chỉnh sửa thông tin cá nhân & Chỉ được sửa họ tên hiển thị và vị trí GPS. Tên đăng nhập và email không được phép thay đổi. & \\
\hline
6 & ND\_QĐ 6 & Quy định gửi ảnh chẩn đoán & Chỉ chấp nhận file JPG hoặc PNG, dung lượng tối đa 10MB. Hệ thống tự động lấy tọa độ GPS để truy xuất dữ liệu thời tiết. & \\
\hline
7 & ND\_QĐ 7 & Quy định nhập thông số môi trường & Không bắt buộc. Nếu nhập thì hệ thống dùng thông số này để điều chỉnh kết quả chẩn đoán. & \\
\hline
8 & ND\_QĐ 8 & Quy định xem kết quả chẩn đoán & Hiển thị tên bệnh, vùng khoanh trên ảnh gốc, độ tin cậy làm tròn và gợi ý dinh dưỡng tương ứng. & \\
\hline
9 & ND\_QĐ 9 & Quy định tra cứu lịch sử & Chỉ hiển thị lịch sử của chính người đang đăng nhập. Lọc theo khoảng thời gian hoặc tên bệnh. Hiển thị 10 dòng mỗi trang. & \\
\hline
10 & ND\_QĐ 10 & Quy định xem chi tiết lịch sử & Hiển thị đầy đủ kết quả chẩn đoán, vùng khoanh và gợi ý dinh dưỡng của lần chẩn đoán đó. & \\
\hline
11 & ND\_QĐ 11 & Quy định gửi phản hồi & Chọn tên bệnh thực tế từ danh sách. Sau khi gửi không được sửa. Phiếu ở trạng thái ``Chờ xử lý''. & \\
\hline
12 & ND\_QĐ 12 & Quy định xem kết quả phản hồi & Hiển thị trạng thái phiếu, tên Admin xử lý và lời nhắn phản hồi. & \\
\hline
\end{longtable}

% ============================================================
\vspace{0.3cm}
\textbf{3. Chi tiết từng Quy định --- Nông dân}

\vspace{0.3cm}
\noindent\textbf{ND\_QĐ 1 --- Quy định Đăng ký tài khoản}

Hệ thống chỉ cho phép tạo tài khoản khi người dùng nhập đầy đủ và hợp lệ các thông tin sau:
\begin{itemize}
    \item \textbf{Tên đăng nhập:} chưa tồn tại trong hệ thống, chỉ gồm chữ cái, chữ số và dấu gạch dưới.
    \item \textbf{Địa chỉ email:} đúng định dạng và chưa được đăng ký trước đó.
    \item \textbf{Mật khẩu:} tối thiểu 8 ký tự, gồm ít nhất 1 chữ hoa, 1 chữ số và 1 ký tự đặc biệt.
    \item \textbf{Xác nhận mật khẩu:} phải gõ lại chính xác giống mật khẩu đã nhập.
\end{itemize}
\textbf{Quy tắc xử lý:} Tạo tài khoản thành công thì tự động chuyển vào màn hình chính.

\vspace{0.3cm}
\noindent\textbf{ND\_QĐ 2 --- Quy định Đăng nhập}

Người dùng cần nhập 2 thông tin sau:
\begin{itemize}
    \item \textbf{Tên đăng nhập hoặc email:} phải tồn tại trong hệ thống.
    \item \textbf{Mật khẩu:} phải khớp với tài khoản tương ứng.
\end{itemize}
\textbf{Quy tắc xử lý:} Nếu sai, hệ thống chỉ thông báo chung ``Tên đăng nhập hoặc mật khẩu không chính xác''. Nhập sai 5 lần liên tiếp thì khóa tài khoản 15 phút.

\vspace{0.3cm}
\noindent\textbf{ND\_QĐ 3 --- Quy định Khôi phục mật khẩu}

Thực hiện theo 4 bước tuần tự:
\begin{enumerate}
    \item Nhập \textbf{địa chỉ email} đã đăng ký trong hệ thống.
    \item Hệ thống gửi \textbf{mã OTP} 6 chữ số về email --- hiệu lực 10 phút, nhập sai 3 lần thì hủy mã.
    \item Nhập \textbf{mật khẩu mới} --- không được trùng mật khẩu cũ, đáp ứng độ mạnh như khi đăng ký.
    \item Nhập \textbf{xác nhận mật khẩu mới} --- phải gõ chính xác giống bước 3.
\end{enumerate}

\vspace{0.3cm}
\noindent\textbf{ND\_QĐ 4 --- Quy định Xem thông tin cá nhân}

Người dùng chỉ xem được thông tin của chính mình sau khi đã đăng nhập. Thông tin hiển thị gồm: tên đăng nhập, họ tên, địa chỉ email, vị trí GPS và ngày tham gia hệ thống.

\vspace{0.3cm}
\noindent\textbf{ND\_QĐ 5 --- Quy định Chỉnh sửa thông tin cá nhân}

Người dùng chỉ được chỉnh sửa các thông tin sau:
\begin{itemize}
    \item \textbf{Họ tên hiển thị:} tên hiển thị trên giao diện.
    \item \textbf{Vị trí GPS:} kéo thả ghim trên bản đồ để cập nhật tọa độ mặc định.
\end{itemize}
\textbf{Quy tắc xử lý:} Tên đăng nhập và địa chỉ email không được phép thay đổi.

\vspace{0.3cm}
\noindent\textbf{ND\_QĐ 6 --- Quy định Gửi ảnh chẩn đoán}

Người dùng cần cung cấp:
\begin{itemize}
    \item \textbf{Ảnh lá lúa:} định dạng JPG hoặc PNG, dung lượng tối đa 10MB.
\end{itemize}
\textbf{Quy tắc xử lý:} Hệ thống tự động lấy tọa độ GPS từ hồ sơ người dùng để truy xuất dữ liệu thời tiết phục vụ chẩn đoán.

\vspace{0.3cm}
\noindent\textbf{ND\_QĐ 7 --- Quy định Nhập thông số môi trường}

Không bắt buộc. Người dùng có thể cung cấp thêm:
\begin{itemize}
    \item \textbf{Vị trí hiện tại:} kéo thả ghim trên bản đồ nếu khác với vị trí trong hồ sơ.
    \item \textbf{Mô tả thực tế:} ghi chú tình trạng ruộng lúa (ví dụ: ngập úng, sương mù).
\end{itemize}
\textbf{Quy tắc xử lý:} Nếu người dùng cung cấp thông tin này, hệ thống sẽ dùng để điều chỉnh kết quả chẩn đoán.

\vspace{0.3cm}
\noindent\textbf{ND\_QĐ 8 --- Quy định Xem kết quả chẩn đoán}

Hệ thống hiển thị kết quả gồm:
\begin{itemize}
    \item \textbf{Tên bệnh:} tên bệnh được chẩn đoán.
    \item \textbf{Vùng bệnh:} khoanh vùng trực tiếp trên ảnh gốc tại vị trí AI phát hiện.
    \item \textbf{Độ tin cậy:} phần trăm độ chính xác của kết quả, được làm tròn.
    \item \textbf{Gợi ý dinh dưỡng:} danh sách hoạt chất và hướng dẫn sử dụng tương ứng với bệnh được chẩn đoán.
\end{itemize}

\vspace{0.3cm}
\noindent\textbf{ND\_QĐ 9 --- Quy định Tra cứu lịch sử chẩn đoán}

Người dùng có thể lọc lịch sử theo:
\begin{itemize}
    \item \textbf{Từ ngày / Đến ngày:} khoảng thời gian cần tra cứu.
    \item \textbf{Tên bệnh:} chọn tên bệnh cụ thể để lọc.
\end{itemize}
\textbf{Quy tắc xử lý:} Chỉ hiển thị lịch sử của chính người đang đăng nhập. Danh sách hiển thị 10 dòng mỗi trang.

\vspace{0.3cm}
\noindent\textbf{ND\_QĐ 10 --- Quy định Xem chi tiết một lần chẩn đoán}

Hệ thống hiển thị đầy đủ thông tin của lần chẩn đoán đó gồm: ảnh gốc, vùng khoanh bệnh, tên bệnh, độ tin cậy và gợi ý dinh dưỡng tương ứng.

\vspace{0.3cm}
\noindent\textbf{ND\_QĐ 11 --- Quy định Gửi phản hồi kết quả}

Người dùng cần cung cấp:
\begin{itemize}
    \item \textbf{Loại bệnh thực tế:} chọn từ danh sách có sẵn. Nếu chọn ``Khác'' thì bắt buộc nhập tên bệnh vào ô trống.
    \item \textbf{Lời nhắn kèm theo:} mô tả thêm lý do hoặc tình trạng thực tế (không bắt buộc).
\end{itemize}
\textbf{Quy tắc xử lý:} Sau khi gửi, phiếu ở trạng thái ``Chờ xử lý'' và không được phép chỉnh sửa.

\vspace{0.3cm}
\noindent\textbf{ND\_QĐ 12 --- Quy định Xem kết quả xử lý phản hồi}

Hệ thống hiển thị thông tin gồm: trạng thái phiếu (Chờ xử lý / Chấp nhận / Từ chối), tên Admin xử lý và lời nhắn phản hồi từ Admin.

% ============================================================
\vspace{0.3cm}
\textbf{4. Các Biểu mẫu --- Nông dân}

\vspace{0.3cm}
\noindent\textbf{ND\_BM 1: PHIẾU ĐĂNG KÝ TÀI KHOẢN}

\noindent Tên đăng nhập: \dotfill \\
Địa chỉ email: \dotfill \\
Mật khẩu: \dotfill \\
Xác nhận mật khẩu: \dotfill

\vspace{0.3cm}
\noindent\textbf{ND\_BM 2: PHIẾU CHỈNH SỬA THÔNG TIN CÁ NHÂN}

\noindent Tên đăng nhập: \dotfill \textit{(không được sửa)} \\
Địa chỉ email: \dotfill \textit{(không được sửa)} \\
Họ tên hiển thị: \dotfill \\
Vị trí GPS: [Bản đồ] \textit{(kéo thả ghim để cập nhật)}

\vspace{0.3cm}
\noindent\textbf{ND\_BM 3: PHIẾU GỬI ẢNH CHẨN ĐOÁN}

\noindent Ảnh lá lúa: [Chọn file] \textit{(JPG/PNG, tối đa 10MB)} \\
Vị trí GPS hiện tại: \dotfill \textit{(tự động lấy từ hồ sơ)} \\
Thông số môi trường bổ sung: [Xem ND\_BM 3b] \textit{(tùy chọn)}

\vspace{0.2cm}
\noindent\textbf{ND\_BM 3b: PHIẾU NHẬP THÔNG SỐ MÔI TRƯỜNG} \textit{(tùy chọn)}

\noindent Vị trí hiện tại: [Bản đồ] \textit{(kéo thả ghim nếu khác hồ sơ)} \\
Mô tả thực tế: \dotfill \\
\dotfill \\
\textit{(Ví dụ: ngập úng, sương mù, ruộng mới phun thuốc...)}

\vspace{0.3cm}
\noindent\textbf{ND\_BM 4: PHIẾU PHẢN HỒI KẾT QUẢ CHẨN ĐOÁN}

\noindent Mã chẩn đoán: \dotfill \quad Ngày gửi: \dotfill \\
AI chẩn đoán: \dotfill \quad Độ tin cậy: \dots\% \\
Loại bệnh thực tế: [Danh sách chọn] \textit{(nếu chọn ``Khác'' thì nhập tên bệnh: \dotfill)} \\
Lời nhắn kèm theo: \dotfill \\
\dotfill

% ============================================================
% BỘ PHẬN: QUẢN TRỊ VIÊN
% ============================================================
\section*{B. BỘ PHẬN: QUẢN TRỊ VIÊN (Mã số: AD)}

\textbf{Mục đích:} Quản lý hệ thống, danh mục bệnh, kho tri thức dinh dưỡng, tài khoản người dùng và theo dõi hiệu suất mô hình AI.

\vspace{0.3cm}
\textbf{1. Bảng yêu cầu chức năng nghiệp vụ}

\begin{longtable}{|c|p{3cm}|p{2cm}|p{2.5cm}|p{2cm}|p{2.5cm}|}
\hline
\rowcolor{headerblue}
\textbf{STT} & \textbf{Công việc} & \textbf{Loại công việc} & \textbf{Quy định/CT liên quan} & \textbf{Biểu mẫu liên quan} & \textbf{Ghi chú} \\
\hline
\endhead
\multicolumn{6}{|l|}{\cellcolor{notegreen}\textbf{I. Quản lý tài khoản quản trị}} \\
\hline
1 & Đăng nhập hệ thống quản trị & Tra cứu & AD\_QĐ 1 & & \\
\hline
2 & Đổi mật khẩu quản trị & Lưu trữ & AD\_QĐ 1 & & \\
\hline
\multicolumn{6}{|l|}{\cellcolor{notegreen}\textbf{II. Quản lý danh mục}} \\
\hline
3 & Quản lý bệnh lúa & Lưu trữ & AD\_QĐ 2 & AD\_BM 1 & \\
\hline
4 & Quản lý tri thức dinh dưỡng & Lưu trữ & AD\_QĐ 3 & AD\_BM 2 & \\
\hline
\multicolumn{6}{|l|}{\cellcolor{notegreen}\textbf{III. Quản lý người dùng}} \\
\hline
5 & Xem danh sách người dùng & Tra cứu & AD\_QĐ 4 & & \\
\hline
6 & Khóa / Mở khóa tài khoản & Lưu trữ & AD\_QĐ 4 & & \\
\hline
\multicolumn{6}{|l|}{\cellcolor{notegreen}\textbf{IV. Quản lý phản hồi}} \\
\hline
7 & Xem danh sách phản hồi & Tra cứu & AD\_QĐ 5 & & \\
\hline
8 & Xử lý phản hồi & Lưu trữ & AD\_QĐ 5 & AD\_BM 3 & \\
\hline
\multicolumn{6}{|l|}{\cellcolor{notegreen}\textbf{V. Hệ thống}} \\
\hline
9 & Cập nhật mô hình AI & Lưu trữ & AD\_QĐ 6 & & \\
\hline
10 & Xem thống kê \& báo cáo & Tra cứu, Kết xuất & AD\_QĐ 7 & AD\_BM 4 & \\
\hline
\end{longtable}

% ============================================================
\vspace{0.3cm}
\textbf{2. Bảng Quy định / Công thức liên quan --- Quản trị viên}

\begin{longtable}{|c|p{1.8cm}|p{3.5cm}|p{6.5cm}|p{1cm}|}
\hline
\rowcolor{headerblue}
\textbf{STT} & \textbf{Mã số} & \textbf{Tên Quy định} & \textbf{Mô tả chi tiết} & \textbf{Ghi chú} \\
\hline
\endhead
1 & AD\_QĐ 1 & Quy định tài khoản quản trị & Đăng nhập bằng tên đăng nhập và mật khẩu. Nhập sai 5 lần thì khóa 30 phút. Đổi mật khẩu phải đáp ứng độ mạnh như khi đăng ký. & \\
\hline
2 & AD\_QĐ 2 & Quy định quản lý bệnh lúa & Thêm, sửa, ẩn bệnh lúa. Tên bệnh phải là duy nhất. Chỉ xóa vĩnh viễn khi chưa có dữ liệu lịch sử. & \\
\hline
3 & AD\_QĐ 3 & Quy định quản lý tri thức dinh dưỡng & Thêm, sửa, xóa thông tin dinh dưỡng. Mỗi bản ghi phải gắn với một loại bệnh cụ thể trong danh mục. & \\
\hline
4 & AD\_QĐ 4 & Quy định quản lý người dùng & Admin chỉ được xem thông tin định danh, không được chỉnh sửa. Chỉ được phép khóa hoặc mở khóa tài khoản. & \\
\hline
5 & AD\_QĐ 5 & Quy định xử lý phản hồi & Xem xét ảnh gốc, kết quả AI và kết quả nông dân báo cáo. Chốt kết quả Chấp nhận hoặc Từ chối và gửi lời nhắn phản hồi. & \\
\hline
6 & AD\_QĐ 6 & Quy định cập nhật mô hình AI & Tải lên file mô hình mới kèm tên phiên bản. Admin dùng công tắc kích hoạt để áp dụng cho toàn hệ thống. & \\
\hline
7 & AD\_QĐ 7 & Quy định thống kê \& báo cáo & Hiển thị bản đồ nhiệt dịch bệnh theo tọa độ GPS và thống kê tỷ lệ chẩn đoán đúng/sai theo từng bệnh. & \\
\hline
\end{longtable}

% ============================================================
\vspace{0.3cm}
\textbf{3. Chi tiết từng Quy định --- Quản trị viên}

\vspace{0.3cm}
\noindent\textbf{AD\_QĐ 1 --- Quy định Tài khoản Quản trị}

Quản trị viên cần nhập 2 thông tin để đăng nhập:
\begin{itemize}
    \item \textbf{Tên đăng nhập:} phải tồn tại trong hệ thống.
    \item \textbf{Mật khẩu:} phải khớp với tài khoản quản trị.
\end{itemize}
\textbf{Quy tắc xử lý:} Nhập sai 5 lần liên tiếp thì khóa tài khoản 30 phút. Khi đổi mật khẩu, mật khẩu mới phải đáp ứng độ mạnh như khi đăng ký và không được trùng mật khẩu cũ.

\vspace{0.3cm}
\noindent\textbf{AD\_QĐ 2 --- Quy định Quản lý bệnh lúa}

Mỗi bản ghi bệnh gồm các thông tin:
\begin{itemize}
    \item \textbf{Tên bệnh:} phải là duy nhất trong hệ thống, dễ hiểu cho nông dân.
    \item \textbf{Tên khoa học:} tên tiếng Anh để đối chiếu chuyên môn.
    \item \textbf{Mô tả triệu chứng:} đặc điểm nhận dạng trên lá hoặc thân lúa.
\end{itemize}
\textbf{Quy tắc xử lý:}
\begin{itemize}
    \item Thêm: tên bệnh không được trùng lặp.
    \item Sửa: không được thay đổi mã định danh hệ thống.
    \item Xóa: chỉ xóa vĩnh viễn khi chưa có dữ liệu lịch sử, ngược lại chỉ được chuyển sang trạng thái ``Ẩn''.
\end{itemize}

\vspace{0.3cm}
\noindent\textbf{AD\_QĐ 3 --- Quy định Quản lý tri thức dinh dưỡng}

Mỗi bản ghi dinh dưỡng gồm các thông tin:
\begin{itemize}
    \item \textbf{Bệnh lúa áp dụng:} chọn từ danh mục bệnh có sẵn, bắt buộc.
    \item \textbf{Hoạt chất bổ sung:} tên vitamin hoặc khoáng chất cần thiết.
    \item \textbf{Hướng dẫn sử dụng:} liều lượng và cách bón/phun chi tiết.
\end{itemize}
\textbf{Quy tắc xử lý:} Dữ liệu cập nhật hoặc xóa có hiệu lực ngay lập tức trong kho tri thức hệ thống.

\vspace{0.3cm}
\noindent\textbf{AD\_QĐ 4 --- Quy định Quản lý người dùng}

Admin chỉ được thực hiện các thao tác sau:
\begin{itemize}
    \item \textbf{Xem thông tin:} tên đăng nhập, email, số lượt chẩn đoán, số lượt phản hồi.
    \item \textbf{Trạng thái tài khoản:} khóa hoặc mở khóa tùy theo vi phạm.
\end{itemize}
\textbf{Quy tắc xử lý:} Tên đăng nhập và email của người dùng tuyệt đối không được chỉnh sửa.

\vspace{0.3cm}
\noindent\textbf{AD\_QĐ 5 --- Quy định Xử lý phản hồi}

Thực hiện theo 4 bước tuần tự:
\begin{enumerate}
    \item Xem danh sách phiếu chưa xử lý, ưu tiên phiếu gửi sớm nhất.
    \item Xem xét song song: ảnh gốc, kết quả AI và kết quả nông dân báo cáo.
    \item Nhập \textbf{lời nhắn phản hồi} gửi lại cho nông dân.
    \item Chốt \textbf{trạng thái xử lý:} Chấp nhận (AI sai) hoặc Từ chối (AI đúng).
\end{enumerate}
\textbf{Quy tắc xử lý:} Nếu chọn ``Chấp nhận'', hệ thống tự động lưu ảnh vào tập dữ liệu để nâng cấp mô hình AI sau này.

\vspace{0.3cm}
\noindent\textbf{AD\_QĐ 6 --- Quy định Cập nhật mô hình AI}

Admin cần cung cấp:
\begin{itemize}
    \item \textbf{File mô hình:} file chứa mô hình AI đã được huấn luyện mới.
    \item \textbf{Tên phiên bản:} tên mô tả đợt nâng cấp.
\end{itemize}
\textbf{Quy tắc xử lý:} Admin dùng công tắc ``Kích hoạt'' để áp dụng mô hình mới cho toàn hệ thống.

\vspace{0.3cm}
\noindent\textbf{AD\_QĐ 7 --- Quy định Thống kê \& Báo cáo}

Hệ thống hiển thị 2 loại báo cáo:
\begin{itemize}
    \item \textbf{Bản đồ dịch bệnh:} hiển thị bản đồ nhiệt tình hình dịch bệnh theo tọa độ GPS của các ca chẩn đoán.
    \item \textbf{Hiệu suất AI:} thống kê tỷ lệ chẩn đoán đúng/sai theo từng loại bệnh dựa trên kết quả xử lý phản hồi.
\end{itemize}

% ============================================================
\vspace{0.3cm}
\textbf{4. Các Biểu mẫu --- Quản trị viên}

\vspace{0.3cm}
\noindent\textbf{AD\_BM 1: PHIẾU QUẢN LÝ BỆNH LÚA}

\noindent Tên bệnh: \dotfill \\
Tên khoa học: \dotfill \\
Mô tả triệu chứng: \dotfill \\
\dotfill \\
Trạng thái: $\square$ Hiển thị \quad $\square$ Ẩn

\vspace{0.3cm}
\noindent\textbf{AD\_BM 2: PHIẾU QUẢN LÝ TRI THỨC DINH DƯỠNG}

\noindent Bệnh lúa áp dụng: [Danh sách chọn] \dotfill \\
Hoạt chất bổ sung: \dotfill \\
Hướng dẫn sử dụng: \dotfill \\
\dotfill

\vspace{0.3cm}
\noindent\textbf{AD\_BM 3: PHIẾU XỬ LÝ PHẢN HỒI}

\noindent Mã phản hồi: \dotfill \quad Người dùng gửi: \dotfill \\
Ảnh gốc: [Hiển thị ảnh lá lúa] \\
AI chẩn đoán: \dotfill \quad Độ tin cậy: \dots\% \\
Nông dân báo cáo: \dotfill \\
Lời nhắn của nông dân: \dotfill \\
Trạng thái xử lý: $\square$ Chấp nhận (AI sai) \quad $\square$ Từ chối (AI đúng) \\
Lời nhắn phản hồi: \dotfill \\
\dotfill

\vspace{0.3cm}
\noindent\textbf{AD\_BM 4: BÁO CÁO HIỆU SUẤT AI}

\noindent Tháng: \dotfill \quad Năm: \dotfill

\begin{longtable}{|c|p{2cm}|p{4cm}|p{3cm}|p{3cm}|}
\hline
\rowcolor{headerblue}
\textbf{STT} & \textbf{Mã bệnh} & \textbf{Tên bệnh} & \textbf{Số lượt nhận diện} & \textbf{Số lượt AI đoán sai} \\
\hline
1 & & & & \\
\hline
2 & & & & \\
\hline
3 & & & & \\
\hline
\end{longtable}
\noindent Ngày lập báo cáo: \dotfill \quad Người lập: \dotfill

\hrulefill

% ============================================================
\subsubsection*{3.2 Xác định yêu cầu chức năng hệ thống và yêu cầu chất lượng}

\textbf{Bảng yêu cầu chức năng hệ thống:}
\begin{longtable}{|c|p{3.5cm}|p{8.5cm}|p{1cm}|}
\hline
\rowcolor{headerblue}
\textbf{STT} & \textbf{Nội dung} & \textbf{Mô tả chi tiết} & \textbf{Ghi chú} \\
\hline
\endhead
1 & Tích hợp thời tiết & Truy xuất dữ liệu thời tiết dựa trên tọa độ thực tế của nông dân để cung cấp bối cảnh môi trường. & \\
\hline
2 & Quản lý tri thức & Hệ thống đồng bộ các hướng dẫn dinh dưỡng để AI có thể tư vấn linh hoạt. & \\
\hline
3 & Bảo mật dữ liệu & Mật khẩu được bảo mật an toàn. Toàn bộ ảnh và dữ liệu thực tế được lưu trữ phục vụ tái huấn luyện. & \\
\hline
4 & Gửi mã xác thực & Tự động gửi mã xác thực qua email cho các luồng xác minh tài khoản. & \\
\hline
\end{longtable}

\textbf{Bảng yêu cầu về chất lượng hệ thống:}
\begin{longtable}{|c|p{3cm}|p{2cm}|p{6cm}|p{1cm}|}
\hline
\rowcolor{headerblue}
\textbf{STT} & \textbf{Nội dung} & \textbf{Tiêu chuẩn} & \textbf{Mô tả chi tiết} & \textbf{Ghi chú} \\
\hline
\endhead
1 & Cập nhật mô hình & Tiến hóa & Có thể thay thế mô hình nhận diện mới mà không làm gián đoạn người dùng. & \\
\hline
2 & Tính tiện dụng & Tiện dụng & Giao diện tối ưu cho thiết bị di động, thông báo rõ ràng bằng tiếng Việt. & \\
\hline
3 & Tốc độ xử lý & Hiệu quả & Phản hồi kết quả chẩn đoán trong vòng 5--7 giây. & \\
\hline
4 & Độ tương thích & Tương thích & Hoạt động ổn định trên các trình duyệt phổ biến của điện thoại thông minh. & \\
\hline
5 & Tính kịp thời & Hiệu quả & Mã xác thực gửi đến người dùng trong vòng 30 giây. & \\
\hline
\end{longtable}

\end{document}
